\newpage\SourceLeaf{64}\ArithmeticHead{57}{R}{KUBISCHE UND BIQUADRATISCHE RESTE. KOMPLEXE GANZE ZAHLEN.}\HistoricalFoot{X\textsubscript{2}.}{8}
\begin{center}\bfseries 23.\end{center}
\noindent\hspace*{\parindent}\textsc{Gauss} hat die Darlegung der Beweise für die biquadratischen Haupttheoreme einer dritten Abhandlung vorbehalten, die aber niemals geschrieben worden ist. Die zweite gibt nur den Wortlaut des allgemeinen quadratischen (W. II, S. 130), sowie auch des allgemeinen biquadratischen Reziprozitätsgesetzes (W. II, S. 138), das zwischen zwei (primären) komplexen Primzahlen besteht, und beweist den sogenannten Ergänzungssatz, durch welchen der biquadratische Charakter der Zahl \(1+i\) und ihrer sogenannten Assoziierten bestimmt wird. Vorauf geht eine, der Arithmetik rationaler Zahlen völlig parallellaufende, nur entsprechend reichere Arithmetik der komplexen Zahlen \(a+bi\), gegründet auf das \textsc{Euklidische} Verfahren zur Bestimmung des grössten gemeinsamen Teilers von zwei solchen Zahlen; hieraus folgt die eindeutige Zerlegbarkeit jeder komplexen Zahl in sogenannte primäre Primfaktoren, deren es ausser der Zahl \(1+i\) zwei Arten gibt, nämlich die reellen Primzahlen \(q\equiv3\) (mod. 4) und die konjugiert komplexen Faktoren der reellen Primzahlen \(p\equiv1\) (mod. 4).

Aus dem verallgemeinerten \textsc{Fermat}schen Lehrsatze, nach welchem, wenn \(\pi\) eine komplexe Primzahl und \(n\) ihre Norm bedeutet, jede durch \(\pi\) nicht teilbare komplexe Zahl \(\kappa\) der Kongruenz
\[
\kappa^{n-1}\equiv1\ (\mathrm{mod.}\ \pi)
\]
genügt, entspringt eine Einteilung aller Reste von \(\pi\) in vier Klassen von gleichviel Zahlen, jenachdem
\[
\kappa^{\frac{n-1}{4}}\equiv1,i,-1,-i\ (\mathrm{mod.}\ \pi)
\]
ist, und die Frage, welcher dieser Klassen eine gegebene Zahl zugehört, ist die Frage nach ihrem biquadratischen Charakter. Verteilt man aber alle Reste in vier Gruppen \(C,C',C'',C'''\) derart, dass aus den Zahlen \(r\) der Gruppe \(C\) die der übrigen durch Multiplikation beziehungsweise mit \(i,-1,-i\) entstehen, ist ferner \(\kappa\) eine durch \(\pi\) nicht teilbare komplexe Zahl und bezeichnet man bezw. mit \(c,c',c'',c'''\) die Anzahlen derjenigen Reste der Produkte \(\kappa r\), die den Gruppen \(C,C',C'',C'''\) angehören, so besteht die Kongruenz
\newpage\SourceLeaf{65}\ArithmeticHead{58}{L}{BACHMANN, ÜBER GAUSS’ ZAHLENTHEORETISCHE ARBEITEN.}
\[
\kappa^{\frac{n-1}{4}}\equiv i^{c'+2c''+3c'''}\ (\mathrm{mod.}\ \pi),
\]
die ersichtlich das Analogon des \textsc{Gauss}schen Lemma ist. Im Verein mit den Eigenschaften des schon in den Comm. Gotting. rec. \textbf{4} (W. II, S. 61) auftretenden Ausdrucks \(\varphi(a,b)\) genügt sie zur Herleitung des biquadratischen Charakters von \(1+i\).

\textsc{Gauss} hat, wie bemerkt, einen Beweis des biquadratischen Fundamentaltheorems nicht veröffentlicht. Umso interessanter ist es, dass in seinem Nachlasse (in Ec 3, abgedruckt W. X1, S. 65) auf vier Oktavseiten geschrieben die flüchtige Skizze eines vollständigen Beweises gefunden worden ist, welcher, auf die Kreisteilung gegründet, dem Beweise VI des quadratischen Reziprozitätsgesetzes entspricht. Es scheint diese Skizze einer späteren Zeit zu entstammen, vielleicht hervorgerufen durch \textsc{Eisensteins} bezügliche Arbeiten\footnote{\textsc{Crelles} Journal für Mathematik \textbf{28}, 1844, S. 53 und 223.}, dessen Symbol \(\left[\frac{\kappa}{m}\right]\) hier auch von \textsc{Gauss} zur Bezeichnung des biquadratischen Charakters von \(\kappa\) bezüglich \(m\) für den Fall einer komplexen Primzahl \(m\) verwandt wird, während er für den allgemeineren Fall einer zusammengesetzten Zahl \(m\) das Symbol \(\left(\frac{\kappa}{m}\right)\) benutzt. Sicherlich wohl enthält sie das, was nach einer von \textsc{Eisenstein} an \textsc{Riemann} getanen und von diesem an \textsc{Schering} weitergegebenen Aussage\begingroup\renewcommand{\thefootnote}{**)}\footnote{Siehe \textsc{Schering}, Göttinger Nachrichten 1879, S. 384.}\endgroup\ \textsc{Gauss} an \textsc{Eisenstein} brieflich mitgeteilt hat. Das Gesetz wird hier in etwas anderer Fassung als in der \textit{Comm. secunda} dahin ausgesprochen, dass, wenn \(m=a+bi,\ M=A+Bi\) zwei ungerade komplexe Zahlen ohne gemeinsamen Teiler und \(a,A\) ungerade sind,
\[
\left(\frac{m}{M}\right)=i^{\frac12 bB}\mathbin{.}\left(\frac{M}{m}\right)
\]
sei. An einer andern Stelle des Nachlasses (in Ec 2) sind noch verschiedene abweichende Formulierungen des Gesetzes vorhanden, deren eine hier mitgeteilt werden möge. Bezeichnen \(m=a+bi,\ m'=a'+b'i\) zwei verschiedene ungerade komplexe Primzahlen, wobei \(a,a'\) ungerade, und setzt man
\[
m'^{\frac14(a^2+b^2-1)}\equiv i^\mu\ (\mathrm{mod.}\ m),\quad m^{\frac14(a'^2+b'^2-1)}\equiv i^{\mu'}\ (\mathrm{mod.}\ m'),
\]
so ist
\[
\mu-\mu'-\tfrac12\bigl((a'-1)b+(a-1)b'+bb'\bigr)\equiv0\ (\mathrm{mod.}\ 4).
\]
\newpage\SourceLeaf{66}\ArithmeticHead{59}{R}{KUBISCHE UND BIQUADRATISCHE RESTE. KOMPLEXE GANZE ZAHLEN.}\HistoricalFoot{}{8*}
\noindent\hspace*{\parindent}Eine im wesentlichen gleiche Formulierung steht auch in dem im September 1813 begonnenen Handbuche 21, Bg, S. 24 (abgedruckt W. X1, S. 55), d. h. hinter den dort ebenfalls niedergeschriebenen Beweisen V und VI des quadratischen Reziprozitätsgesetzes (W. X1, S. 33 und 28), sie ist also vermutlich auch jüngeren Datums als 12. November 1813, was sich mit \textsc{Gauss}’ Angaben in dem erwähnten Briefe an \textsc{Dirichlet} wohl verträgt.
